\chapter{ТЕОРЕТИЧЕСКИЕ ПРЕДПОСЫЛКИ, МОДЕЛИ И МЕТОДОЛОГИЯ}
\label{ch:model_algorithm}

\section{Формальная модель конвейера обработки}
\label{sec:formal_pipeline}

Сквозной маршрут платформы графически представлен на рисунке~\ref{fig:ch2_pipeline}: от приема RTSP-потока до фиксации события в БД и его асинхронной доставки во внешний контур. Эта схема используется как опорная при формализации модели и при описании реализации в главе~\ref{ch:implementation}.

\begin{figure}[H]
\centering
\begin{tikzpicture}[
    node distance=6mm and 7mm,
    every node/.style={font=\small},
    stage/.style={
        draw=cBorder, line width=0.4pt, rounded corners=2pt,
        fill=cBox, text=cText,
        minimum width=22mm, minimum height=11mm,
        align=center, inner sep=3pt
    },
    sink/.style={stage, fill=cBoxAcc},
    src/.style={stage, fill=cBoxAlt},
    arr/.style={-{Stealth[length=2.2mm]}, draw=cArrow, line width=0.5pt}
]
\node[src] (cam) {RTSP\\поток};
\node[stage, right=of cam] (det) {Детектор\\YOLO};
\node[stage, right=of det] (trk) {Сопровождение\\DeepSORT};
\node[stage, right=of trk] (rule) {Правиловой\\движок};
\node[stage, below=10mm of rule] (out) {\texttt{event\_outbox}\\(PostgreSQL)};
\node[stage, left=of out] (br) {Брокер\\RabbitMQ};
\node[sink, left=of br] (wh) {HTTP-\\получатель};
\draw[arr] (cam) -- (det);
\draw[arr] (det) -- node[above,font=\tiny]{bbox} (trk);
\draw[arr] (trk) -- node[above,font=\tiny]{траект.} (rule);
\draw[arr] (rule) -- node[right,font=\tiny]{событие} (out);
\draw[arr] (out) -- node[above,font=\tiny]{публ.} (br);
\draw[arr] (br) -- node[above,font=\tiny]{дост.} (wh);
\end{tikzpicture}
\caption{Сквозной конвейер обработки видеопотока и доставки событий}
\label{fig:ch2_pipeline}
\end{figure}

Платформа рассматривается как композиция взаимосвязанных функций
\begin{equation}
\Phi = \Phi_{\text{outbox}} \circ \Phi_{\text{rules}} \circ \Phi_{\text{track}} \circ \Phi_{\text{detect}} \circ \Phi_{\text{capture}},
\end{equation}
\where
\whereitem{\(i \in \{1, \dots, N\}\)}{индекс камеры}
\whereitemlast{\(\Phi\)}{сквозной конвейер обработки видеопотока для камеры \(i\)} \cite{fastapi_docs,pyav_docs,ultralytics_docs}

Для камеры \(i\) в момент времени \(t\) на вход поступает поток кадров \(x^i_t\). Стадия захвата и нормализации формирует последовательность
\begin{equation}
\hat{x}^i_t = \Phi_{\text{capture}}(x^i_t, m_i),
\end{equation}
\where
\whereitemlast{\(m_i\)}{метаданные источника: временная зона, целевой FPS и стратегия переподключения}

На критическом пути применяется политика замещения последним актуальным кадром: очередь между захватом и инференсом ограничена размером 1. Формально буфер \(\mathcal{Q}_i\) в момент \(t\) определяется как
\begin{equation}
\mathcal{Q}_i(t) =
\begin{cases}
\{\hat{x}^i_t\}, & \text{если буфер пуст;}\\
\{\hat{x}^i_t\}, & \text{если буфер содержит устаревший кадр.}
\end{cases}
\end{equation}
Такой режим эквивалентен оператору замещения последнего элемента и ограничивает накопление задержки \cite{ffmpeg_docs}.

Детектор формирует множество кандидатов
\begin{equation}
\mathcal{D}^i_t = \Phi_{\text{detect}}(\hat{x}^i_t;\theta_d),
\end{equation}
\where
\whereitemlast{\(\theta_d\)}{параметры нейросетевой модели класса YOLO} \cite{redmon2016yolo,redmon2018yolov3,ultralytics_docs}

Модуль сопровождения преобразует детекции в траектории:
\begin{equation}
\mathcal{T}^i_t = \Phi_{\text{track}}(\mathcal{D}^i_t, \mathcal{T}^i_{t-1}; \theta_t).
\end{equation}
Результатом являются активные треки с устойчивыми \texttt{track\_id} и обновленными состояниями \cite{bewley2016sort,wojke2017deepsort}.

Правиловой движок вычисляет предикаты событий:
\begin{equation}
\mathcal{E}^i_t = \Phi_{\text{rules}}(\mathcal{T}^i_t, \mathcal{Z}_i, \mathcal{R}_i),
\end{equation}
\where
\whereitem{\(\mathcal{Z}_i\)}{геометрия зон и линий контроля камеры \(i\)}
\whereitemlast{\(\mathcal{R}_i\)}{конфигурация правил событийного анализа камеры \(i\)}

Наконец, выход конвейера фиксируется транзакционно:
\begin{equation}
(\texttt{events}, \texttt{event\_outbox}) = \Phi_{\text{outbox}}(\mathcal{E}^i_t).
\end{equation}
Тем самым обеспечивается логическая атомарность между фактом обнаруженного события и фактом его последующей публикации во внешний брокер \cite{transactional_outbox,rabbitmq_docs}.

\section{Обоснование архитектуры и технологического стека}
\label{sec:stack_rationale}

В данном разделе формализуется выбор архитектурного решения и технологического стека платформы. Согласно методическим рекомендациям, во второй главе ВКР должна быть не только приведена архитектура и используемый набор технологий, но и обосновано, почему выбраны именно они. Обоснование строится на трех уровнях: архитектурный паттерн, алгоритмическое ядро и инфраструктурный стек.

\subsection{Архитектурный паттерн}

Платформа реализована в виде набора слабо связанных микросервисов, взаимодействующих через явно описанные программные контракты. Сравнение основных архитектурных альтернатив для рассматриваемого класса задач приведено в таблице~\ref{tab:arch_pattern_choice}.

\begin{table}[H]
\caption{Сравнение архитектурных паттернов}
\label{tab:arch_pattern_choice}
\begin{tabularx}{\textwidth}{>{\raggedright\arraybackslash}p{0.27\textwidth}X}
\toprule
Паттерн & Оценка применимости \\
\midrule
Монолитное приложение & Простое развертывание, но смешение алгоритмического и интеграционного контуров повышает риск каскадных отказов и затрудняет независимое масштабирование \\
Модульный монолит & Снижает связность по сравнению с обычным монолитом, однако сохраняет общую модель отказов и общий контур исполнения \\
Микросервисы & Изолируют контуры обработки видео, конфигурации и интеграции; обеспечивают независимое масштабирование, но требуют надежного транспорта сообщений и наблюдаемости \\
Serverless / FaaS & Пригодны для эпизодических задач, но плохо согласуются с непрерывной потоковой обработкой видеоданных \\
\bottomrule
\end{tabularx}
\end{table}

Для платформы видеоаналитики реального времени выбран микросервисный архитектурный шаблон с разделением на контур управления (управление, конфигурация, доступ) и контур обработки данных (обработка видео, генерация и доставка событий). Такое разделение соответствует рекомендациям инженерной литературы по построению событийно-ориентированных систем~\cite{richardson_microservices,transactional_outbox} и обеспечивает выполнение архитектурных факторов, сформулированных в главе 1: отказ интеграционного контура не должен блокировать аналитическое ядро, а замена детектора или модуля сопровождения должна быть локальным изменением.

\subsection{Алгоритмическое ядро: детекция и сопровождение}

Выбор алгоритмов обнаружения и сопровождения объектов рассмотрен в разделе~\ref{sec:related_work} главы~1. Здесь формализуем критерии и итоговое решение в табличной форме (таблица~\ref{tab:algo_choice}).

\begin{table}[H]
\caption{Обоснование выбора алгоритмов обнаружения и сопровождения}
\label{tab:algo_choice}
\begin{tabularx}{\textwidth}{@{}>{\raggedright\arraybackslash}p{0.21\textwidth} >{\raggedright\arraybackslash}p{0.23\textwidth} X@{}}
\toprule
Подсистема & Выбранный метод & Ключевые причины выбора \\
\midrule
Детектор объектов & YOLO-семейство (одноэтапный)~\cite{redmon2016yolo,redmon2018yolov3,ultralytics_docs} & Высокая скорость инференса, многомасштабное представление признаков, доступность открытых реализаций, возможность дообучения на предметных данных, в том числе для малых объектов (БПЛА) \\
Сопровождение объектов & DeepSORT-подход на основе фильтра Калмана и признаков внешнего вида~\cite{wojke2017deepsort,kalman1960,kuhn1955hungarian} & Сочетание кинематического прогноза и признаков внешнего вида снижает число переключений идентификаторов (ID-switches) при перекрытиях; формальная процедура ассоциации совместима с гейтингом и многокамерной изоляцией \\
Событийный анализ & Правиловой движок над траекториями~\cite{transactional_outbox} & Интерпретируемость событий, простота адаптации к локальным регламентам, независимость от конкретной реализации детектора и модуля сопровождения \\
\bottomrule
\end{tabularx}
\end{table}

Альтернативные подходы (двухэтапные и трансформерные детекторы, end-to-end событийные нейросети) обладают преимуществами по точности на сложных сценах, однако предъявляют существенно более высокие требования к вычислительным ресурсам и объему размеченных данных по событиям, что не согласуется с целями практико-ориентированного исследовательского стенда.

\subsection{Веб-фреймворк управляющего сервиса}

Сервис \texttt{api} реализует управляющий контур: аутентификацию, авторизацию, управление конфигурацией камер, зон, правил и доступом к журналу событий. Сравнение основных альтернатив для Python-экосистемы приведено в таблице~\ref{tab:framework_choice}.

\begin{table}[H]
\caption{Сравнение веб-фреймворков для управляющего сервиса}
\label{tab:framework_choice}
\begin{tabularx}{\textwidth}{>{\raggedright\arraybackslash}p{0.18\textwidth}X}
\toprule
Фреймворк & Оценка применимости \\
\midrule
Django & Полнофункциональный фреймворк с ORM и admin-панелью, но избыточен для тонкого API-контура и ограничен по асинхронной модели исполнения \\
Flask & Минималистичный фреймворк, но требует ручной реализации валидации, асинхронности и автодокументации API \\
FastAPI~\cite{fastapi_docs} & Современный асинхронный фреймворк, основанный на Starlette и Pydantic; обеспечивает автоматическую генерацию OpenAPI-схемы, строгую валидацию входных данных и хорошую совместимость с асинхронным стеком \\
\bottomrule
\end{tabularx}
\end{table}

Для управляющего сервиса выбран FastAPI~\cite{fastapi_docs} в сочетании с SQLAlchemy и Alembic~\cite{sqlalchemy_docs,alembic_docs}. Основные мотивы: автоматическая генерация спецификации OpenAPI и контроль контракта на этапе разработки; нативная асинхронная модель исполнения, согласованная с типичной нагрузкой управляющего API; типобезопасные модели данных на основе Pydantic, что снижает риск рассинхронизации схемы и кода.

\subsection{Система управления базами данных}

Хранилище платформы обслуживает структурированные сущности (пользователи, роли, камеры, зоны, правила, события) и таблицу \texttt{event\_outbox}, требующую транзакционной согласованности с записью событий. Сравнение СУБД приведено в таблице~\ref{tab:db_choice}.

\begin{table}[H]
\caption{Сравнение систем управления базами данных}
\label{tab:db_choice}
\begin{tabularx}{\textwidth}{>{\raggedright\arraybackslash}p{0.18\textwidth}X}
\toprule
СУБД & Оценка применимости \\
\midrule
SQLite & Подходит для прототипов, но не предназначена для конкурентной нагрузки и работы под несколькими сервисами \\
\shortstack[l]{MySQL/\\MariaDB} & Зрелая реляционная СУБД, однако ряд особенностей реализации (например, поведение блокировок, поддержка JSON и расширений) уступает PostgreSQL \\
PostgreSQL & Реляционная СУБД с богатой поддержкой расширений, типов JSON/JSONB и продвинутых индексов; обеспечивает надежную транзакционную семантику, необходимую для шаблона транзакционной таблицы исходящих сообщений \\
MongoDB & Документная СУБД, удобная для гибких схем, но менее естественна для строго транзакционной фиксации \texttt{events + event\_outbox} в одной операции \\
\bottomrule
\end{tabularx}
\end{table}

Выбор PostgreSQL~\cite{postgresql_docs} обусловлен необходимостью атомарной фиксации записей в \texttt{events} и \texttt{event\_outbox} в рамках одной транзакции, а также удобством работы с полем полезной нагрузки \texttt{payload} в формате JSONB и наложением уникального ограничения на поле \texttt{dedup\_key}. Эти возможности непосредственно используются в формальной модели надежной доставки.

\subsection{Брокер сообщений}

Интеграционный контур платформы строится на основе шаблона транзакционной таблицы исходящих сообщений с асинхронной публикацией доменных событий во внешние системы. В таблице~\ref{tab:broker_choice} приведено сопоставление основных вариантов.

\begin{table}[H]
\caption{Сравнение брокеров сообщений}
\label{tab:broker_choice}
\begin{tabularx}{\textwidth}{>{\raggedright\arraybackslash}p{0.22\textwidth}X}
\toprule
Брокер & Оценка применимости \\
\midrule
Apache Kafka & Высокая пропускная способность и упорядоченные журналы, но избыточная инфраструктурная сложность для пилотного стенда с числом камер 1--10 \\
Redis Streams & Простота развертывания, но менее зрелые гарантии надежной доставки и менее богатая модель маршрутизации \\
RabbitMQ~\cite{rabbitmq_docs,rabbitmq_confirms,rabbitmq_dlq} & Зрелая модель маршрутизации через обменники, очереди и привязки, подтверждения публикации RabbitMQ и обменник недоставленных сообщений \\
\bottomrule
\end{tabularx}
\end{table}

Выбор RabbitMQ обусловлен сочетанием подтвержденной публикации, штатной поддержки очередей недоставленных сообщений и достаточной для пилотного стенда пропускной способности~\cite{rabbitmq_docs,rabbitmq_confirms,rabbitmq_dlq}. Это позволяет реализовать формализованную в дальнейшем модель повторных попыток и приближенную оценку вероятности успешной доставки без введения избыточно сложной инфраструктуры.

\subsection{Подсистема видеозахвата}

Для приема RTSP-потоков и их декодирования сопоставлены три варианта: \texttt{OpenCV VideoCapture}, \texttt{GStreamer} и \texttt{FFmpeg/PyAV} (см. таблицу~\ref{tab:capture_alternatives} в главе~1). В платформе использован \texttt{FFmpeg/PyAV}~\cite{ffmpeg_docs,pyav_docs} как компромисс между управляемостью потока, зрелостью экосистемы и интеграцией с Python-контуром исполнения. Это решение позволяет реализовать политику замещения кадра последним актуальным кадром, формализованную в разделе~\ref{sec:formal_pipeline}.

\subsection{Наблюдаемость и оркестрация}

Для воспроизводимого развертывания и эксплуатационной приемки выбраны Docker Compose~\cite{docker_docs}, Prometheus~\cite{prometheus_docs} и Grafana~\cite{grafana_docs}. Docker Compose обеспечивает декларативное описание стенда и согласованность окружений между разработкой и приемочными прогонами; Prometheus реализует сбор и хранение временных рядов, Grafana - их визуализацию и настройку предупреждений. Альтернативы (Kubernetes, Datadog, ELK-стек) обладают большей мощностью, но для пилотного односерверного стенда вносят избыточную сложность сопровождения, не оправданную целями исследования.

\subsection{Сводный технологический стек}

Итоговый технологический стек платформы зафиксирован в таблице~\ref{tab:stack_summary}. Для каждой подсистемы указаны выбранное решение и ключевые источники, обосновывающие этот выбор.

\begin{table}[H]
\caption{Сводный технологический стек платформы}
\label{tab:stack_summary}
\begin{tabularx}{\textwidth}{>{\raggedright\arraybackslash}p{0.22\textwidth} >{\raggedright\arraybackslash}p{0.22\textwidth} X}
\toprule
Подсистема & Выбранное решение & Назначение и обоснование \\
\midrule
Управляющий API & FastAPI + SQLAlchemy + Alembic~\cite{fastapi_docs,sqlalchemy_docs,alembic_docs} & Тонкий асинхронный API с автогенерацией OpenAPI-схемы и строгой валидацией данных \\
Хранилище данных & PostgreSQL~\cite{postgresql_docs} & Транзакционная фиксация \texttt{events + event\_outbox}, поддержка JSONB, уникальные индексы для дедупликации \\
Брокер сообщений & RabbitMQ~\cite{rabbitmq_docs,rabbitmq_confirms,rabbitmq_dlq} & Подтвержденная публикация, маршрутизация по типам событий, штатная очередь недоставленных сообщений \\
Видеозахват & FFmpeg и PyAV~\cite{ffmpeg_docs,pyav_docs} & Управляемый декодер RTSP с предсказуемым поведением при сбоях \\
Детектор & YOLO-семейство~\cite{redmon2016yolo,redmon2018yolov3,ultralytics_docs} & Одноэтапная архитектура с приемлемой задержкой инференса \\
Модуль сопровождения & DeepSORT-подход~\cite{wojke2017deepsort,kalman1960,kuhn1955hungarian} & Сочетание кинематического прогноза и признаков внешнего вида \\
Метрики и мониторинг & Prometheus + Grafana~\cite{prometheus_docs,grafana_docs} & Эксплуатационная наблюдаемость и регламент приемки \\
Оркестрация стенда & Docker Compose~\cite{docker_docs} & Декларативный воспроизводимый стенд для разработки и приемки \\
\bottomrule
\end{tabularx}
\end{table}

Сводный стек согласован между собой: все компоненты имеют открытые реализации, поддерживают типовые контракты обмена данными (HTTP, AMQP, SQL) и обеспечивают совместное выполнение требований к надежности доставки, наблюдаемости и воспроизводимости стенда. Тем самым обеспечивается прямая связь между сформулированными в главе 1 требованиями и инженерным контуром, формализуемым ниже.

Композиция сервисов и их связи показаны на рисунке~\ref{fig:ch2_c4_container} в нотации C4 (контейнерный уровень). Слева расположены внешние участники взаимодействия (камеры, оператор), в центре - управляющий контур (\texttt{api}, \texttt{ui}, PostgreSQL), справа - контур обработки событий (\texttt{analytics-worker}, RabbitMQ, \texttt{integration-worker}) и внешние получатели событий. Подсистема наблюдаемости (Prometheus и Grafana) собирает метрики со всех сервисов через интерфейсы \texttt{/metrics}.

\begin{figure}[H]
\centering
\resizebox{\textwidth}{!}{%
\begin{tikzpicture}[
    every node/.style={font=\small},
    cont/.style={
        draw=cBorder, line width=0.4pt, rounded corners=2pt,
        fill=cBox, text=cText,
        minimum width=28mm, minimum height=11mm, align=center, inner sep=3pt
    },
    db/.style={cont, fill=cBoxAcc},
    bus/.style={cont, fill=cBoxAlt},
    obs/.style={cont, fill=cBoxAcc},
    actor/.style={
        draw=cBorder, line width=0.4pt, rounded corners=2pt,
        fill=cBoxAlt, text=cText,
        minimum width=20mm, minimum height=10mm, align=center, inner sep=3pt
    },
    arr/.style={-{Stealth[length=2.2mm]}, draw=cArrow, line width=0.5pt},
    arrd/.style={-{Stealth[length=2.2mm]}, draw=cArrow, line width=0.4pt, dashed}
]
\node[actor] (cam) at (-6.5,2.2) {IP-камеры};
\node[actor] (op)  at (-6.5,0.4) {Оператор};
\node[actor] (ext) at (7.4,-1.6) {Внешние\\системы};
\node[cont] (aw)  at (-3,2.2) {analytics-worker\\\footnotesize детекция, треки, правила};
\node[cont] (ui)  at (-3,0.4) {ui\\\footnotesize React + nginx};
\node[cont] (api) at ( 0.5,0.4) {api\\\footnotesize FastAPI, JWT, RBAC};
\node[db]   (pg)  at ( 0.5,-1.6) {PostgreSQL\\\footnotesize events, event\_outbox};
\node[bus]  (mq)  at ( 4,0.4) {RabbitMQ};
\node[cont] (iw)  at ( 4,-1.6) {integration-worker\\\footnotesize публикация, повторы\\\footnotesize очередь\\\footnotesize недоставленных};
\node[obs]  (prom) at (-3,-1.6) {Prometheus\\\footnotesize + Grafana};
\draw[arr] (cam) -- node[above,font=\scriptsize]{RTSP} (aw);
\draw[arr] (op)  -- (ui);
\draw[arr] (ui)  -- node[above,font=\scriptsize]{REST} (api);
\draw[arr] (api) -- node[right,font=\scriptsize]{SQL} (pg);
\draw[arr] (aw) -| node[pos=0.25,above,font=\scriptsize]{outbox} (pg);
\draw[arr] (pg) -- node[above,font=\scriptsize]{poll} (iw);
\draw[arr] (iw) -- node[right,font=\scriptsize]{publish} (mq);
\draw[arr] (mq.east) -- ++(0.6,0) |- node[pos=0.25,right,font=\scriptsize]{deliver} (ext);
\draw[arrd] (aw) -- node[left,font=\scriptsize]{/metrics} (prom);
\draw[arrd] (api) -- (prom);
\draw[arrd] (iw) -- (prom);
\end{tikzpicture}
}
\caption{Архитектура платформы (C4, контейнерный уровень)}
\label{fig:ch2_c4_container}
\end{figure}

Диаграмма последовательности надежной асинхронной доставки события представлена на рисунке~\ref{fig:ch2_seq_outbox}. Транзакционная фиксация \texttt{events + event\_outbox} в одной операции БД обеспечивает атомарность факта возникновения события; последующая публикация в RabbitMQ и доставка HTTP-получателю выполняются асинхронно с возможностью повторов и перевода в очередь недоставленных сообщений при исчерпании попыток.

\begin{figure}[H]
\centering
\begin{tikzpicture}[
    every node/.style={font=\footnotesize},
    actor/.style={
        draw=cBorder, line width=0.4pt, rounded corners=2pt,
        fill=cBoxAlt, text=cText,
        minimum width=22mm, minimum height=8mm, align=center
    },
    msg/.style={-{Stealth[length=2mm]}, draw=cArrow, line width=0.4pt},
    note/.style={font=\scriptsize\itshape, text=cText, align=center}
]
\node[actor] (aw)  at (0,0)   {analytics-worker};
\node[actor] (db)  at (3.5,0) {PostgreSQL};
\node[actor] (iw)  at (7,0)   {integration-worker};
\node[actor] (mq)  at (10.5,0){RabbitMQ};
\node[actor] (wh)  at (14,0)  {HTTP-\\получатель};
\foreach \n in {aw,db,iw,mq,wh} {\draw[gray,dashed,line width=0.3pt] (\n.south) -- ++(0,-7);}

\draw[msg] ($(aw.south)+(0,-0.5)$) -- node[above,font=\scriptsize]{INSERT events + event\_outbox (TX)} ($(db.south)+(0,-0.5)$);
\node[note] at (1.7,-1.1) {фиксация};

\draw[msg] ($(iw.south)+(0,-1.6)$) -- node[above,font=\scriptsize]{SELECT pending} ($(db.south)+(0,-1.6)$);
\draw[msg] ($(db.south)+(0,-2.1)$) -- node[above,font=\scriptsize]{записи} ($(iw.south)+(0,-2.1)$);

\draw[msg] ($(iw.south)+(0,-2.8)$) -- node[above,font=\scriptsize]{publish(event)} ($(mq.south)+(0,-2.8)$);
\node[note] at (8.7,-3.3) {опубликовано};

\draw[msg] ($(mq.south)+(0,-4)$) -- node[above,font=\scriptsize]{POST /webhook} ($(wh.south)+(0,-4)$);
\draw[msg] ($(wh.south)+(0,-4.5)$) -- node[above,font=\scriptsize]{2xx} ($(mq.south)+(0,-4.5)$);

\draw[msg] ($(iw.south)+(0,-5.4)$) -- node[above,font=\scriptsize]{UPDATE outbox set published\_at} ($(db.south)+(0,-5.4)$);

\draw[msg,dashed] ($(wh.south)+(0,-6.2)$) -- node[above,font=\scriptsize,align=center]{ошибка $\rightarrow$ повтор с задержкой\\затем очередь недоставленных} ($(mq.south)+(0,-6.2)$);
\end{tikzpicture}
\caption{Диаграмма последовательности надежной доставки события через транзакционную таблицу исходящих сообщений}
\label{fig:ch2_seq_outbox}
\end{figure}

\subsection{Логическая модель работы платформы}
\label{subsec:logic_model}

Логика работы платформы укрупненно описывается следующей последовательностью шагов:
\begin{enumerate}
    \item \textit{Управляющий контур.} Администратор системы через API регистрирует камеры, описывает геометрию зон и линий, задает набор правил событийного анализа и распределяет роли пользователей;
    \item \textit{Активация камеры.} Сервис \texttt{analytics-worker} получает конфигурацию из API и инициирует прием RTSP-потока, политику декодирования и работу детектора;
    \item \textit{Основной цикл обработки.} Для каждого актуального кадра выполняется последовательность «детекция $\rightarrow$ фильтрация классов $\rightarrow$ обновление модуля сопровождения $\rightarrow$ применение правил»;
    \item \textit{Транзакционная фиксация.} При выполнении предиката правила в одной транзакции записываются строки в \texttt{events} и \texttt{event\_outbox}, что исключает потерю события на границе БД;
    \item \textit{Асинхронная публикация.} Сервис \texttt{integration-worker} опрашивает \texttt{event\_outbox}, публикует сообщения в RabbitMQ через подтвержденный канал и обновляет статус записей;
    \item \textit{Доставка во внешний контур.} Тот же сервис принимает сообщения из очереди \texttt{integration.webhook.v1}, выполняет HTTP-доставку внешнему получателю и фиксирует попытки в \texttt{delivery\_attempts};
    \item \textit{Контроль и наблюдаемость.} Все сервисы экспортируют метрики Prometheus и журналы, на основе которых рассчитываются эксплуатационные показатели платформы;
    \item \textit{Операторские сценарии.} Сервис пользовательского интерфейса предоставляет роль-зависимый интерфейс для мониторинга камер, журнала событий и аналитических отчетов.
\end{enumerate}

Описанная логика непосредственно соответствует формальному определению конвейера~$\Phi$ из раздела~\ref{sec:formal_pipeline} и связывает обоснованный технологический стек с математической моделью обработки.

\section{Модель детекции объектов и фильтрации классов}

\subsection{Выход детектора и базовые обозначения}

Для кадра \(\hat{x}_t\) детектор возвращает множество боксов
\begin{equation}
\mathcal{D}_t = \{d_k\}_{k=1}^{K_t}, \quad
d_k = (b_k, c_k, p_k),
\end{equation}
\where
\whereitem{\(b_k = (x_k, y_k, w_k, h_k)\)}{прямоугольник локализации объекта}
\whereitem{\(c_k\)}{класс обнаруженного объекта}
\whereitemlast{\(p_k \in [0,1]\)}{оценка достоверности детекции}

Для подавления шумовых кандидатов используется порог
\begin{equation}
\mathcal{D}^{(\tau)}_t = \{d_k \in \mathcal{D}_t \mid p_k \ge \tau_{\text{conf}}\}.
\end{equation}
После пороговой фильтрации применяется NMS по IoU:
\begin{equation}
\operatorname{IoU}(b_a,b_b) = \frac{|b_a \cap b_b|}{|b_a \cup b_b|},
\end{equation}
и кандидаты с \(\operatorname{IoU} > \tau_{\text{iou}}\) подавляются в пользу бокса с большим \(p_k\) \cite{redmon2018yolov3,coco_dataset}.

\subsection{Фильтрация целевых классов}

Множество целевых классов не является жестко зашитой константой алгоритма и задается профилем применения платформы:
\begin{equation}
\mathcal{C}_{\text{target}}^{(p)} \subseteq \mathcal{C}_{\text{model}},
\end{equation}
\where
\whereitem{\(p\)}{профиль применения платформы}
\whereitemlast{\(\mathcal{C}_{\text{model}}\)}{множество классов, поддерживаемых подключенной моделью детекции}

Для базового охранного профиля используется множество
\begin{equation}
\mathcal{C}_{\text{target}}^{(\text{security})} =
\{\texttt{person}, \texttt{vehicle}\},
\end{equation}
\where
\whereitemlast{\texttt{vehicle}}{агрегированный класс транспорта (\texttt{car}, \texttt{bus}, \texttt{truck}, \texttt{motorcycle}) на уровне API-конфигурации}
Для демонстрационного профиля обнаружения БПЛА используется другое множество целевых классов:
\begin{equation}
\mathcal{C}_{\text{target}}^{(\text{uav})} =
\{\texttt{bird}, \texttt{drone}\}.
\end{equation}

Итоговое множество детекций, передаваемых в модуль сопровождения для выбранного профиля \(p\):
\begin{equation}
\mathcal{D}^{*}_t(p) =
\{d_k \in \mathcal{D}^{(\tau)}_t \mid c_k \in \mathcal{C}_{\text{target}}^{(p)}\}.
\end{equation}
Выбор CPU-first профиля для разработки означает, что \(\tau_{\text{conf}}\), входное разрешение и частота инференса подбираются так, чтобы сохранять ограничение по задержке в целевых сценариях \cite{ultralytics_docs}.

\subsection{Оценка качества детекции}

Для оффлайн-оценки параметров детектора вводится стандартная метрика среднего качества локализации:
\begin{equation}
\operatorname{mAP} = \frac{1}{|\mathcal{C}_{\text{target}}|}\sum_{c \in \mathcal{C}_{\text{target}}}\operatorname{AP}_c,
\end{equation}
\where
\whereitemlast{\(\operatorname{AP}_c\)}{площадь под кривой precision-recall для класса \(c\) при заданных порогах IoU}

Однако в рамках платформы реального времени метрики детекции анализируются совместно с показателями последующих этапов: устойчивостью траекторий и точностью генерации событий. Это исключает ситуацию, когда рост mAP не приводит к улучшению эксплуатационного качества контура.

\section{Модель сопровождения объектов}

\subsection{Состояние трека и прогноз}

Состояние трека \(j\) в модели DeepSORT задается вектором
\begin{equation}
\mathbf{s}_{j,t} = [u, v, \gamma, h, \dot{u}, \dot{v}, \dot{\gamma}, \dot{h}]^\top,
\end{equation}
\where
\whereitem{\((u,v)\)}{центр бокса}
\whereitem{\(\gamma\)}{отношение сторон бокса}
\whereitem{\(h\)}{высота бокса}
\whereitemlast{\(\dot{u}, \dot{v}, \dot{\gamma}, \dot{h}\)}{скорости изменения соответствующих компонент} \cite{wojke2017deepsort}

Прогноз состояния выполняется линейной динамической моделью (фильтр Калмана):
\begin{equation}
\mathbf{s}_{j,t|t-1} = \mathbf{F}\mathbf{s}_{j,t-1|t-1} + \mathbf{w}_t,
\end{equation}
\where
\whereitemlast{\(\mathbf{w}_t\)}{шум процесса}

\subsection{Ассоциация детекций и траекторий}

Для пары «траектория--детекция» рассчитывается композиционная стоимость:
\begin{equation}
\mathcal{L}(j,k) = \lambda \cdot d_{\text{motion}}(j,k) + (1-\lambda)\cdot d_{\text{app}}(j,k),
\end{equation}
\where
\whereitem{\(d_{\text{motion}}\)}{расстояние Махаланобиса относительно прогноза}
\whereitemlast{\(d_{\text{app}}\)}{косинусная дистанция признаков внешнего вида}

Оптимальное сопоставление выполняется венгерским алгоритмом \cite{kuhn1955hungarian}. Пары с превышением порогов по \(d_{\text{motion}}\) или \(d_{\text{app}}\) отбрасываются.

\subsection{Устойчивость идентификаторов и fallback-логика}

Практическая устойчивость модуля сопровождения оценивается метриками \(ID\text{-}switches\), \(MOTA\), \(MOTP\) и долей прерванных траекторий \cite{bernardin2008motmetrics,motchallenge}.

Для повышения устойчивости в контуре исполнения применяется fallback-логика:
\begin{enumerate}
    \item при кратковременной потере детекции траектория удерживается в состоянии \texttt{tentative/lost} до таймаута \(T_{\text{lost}}\);
    \item при превышении \(T_{\text{lost}}\) траектория закрывается, но ее финальные атрибуты сохраняются для правил \texttt{dwell\_time};
    \item при деградации appearance-эмбеддингов вес \(\lambda\) динамически смещается в сторону кинематической составляющей.
\end{enumerate}

Изоляция состояния по камерам формально задается как
\begin{equation}
\mathcal{T}_t = \bigsqcup_{i=1}^{N}\mathcal{T}^i_t,
\end{equation}
то есть объединение траекторий без пересечения пространств идентификаторов между различными камерами.

\section{Правиловая модель событий и дедупликация}

\subsection{Событие как предикат над траекторией}

Каждое правило \(r \in \mathcal{R}_i\) задает предикат
\begin{equation}
P_r(\tau_j, \mathcal{Z}_i, t) \in \{0,1\}.
\end{equation}
Если \(P_r = 1\), создается событие
\begin{equation}
e =
\begin{aligned}[t]
(&\texttt{event\_id},\ \texttt{event\_type},\ \texttt{camera\_id},\ \texttt{track\_id},\\
 &\texttt{occurred\_at},\ \texttt{severity},\ \texttt{dedup\_key},\ \dots).
\end{aligned}
\end{equation}
\subsection{Правило \texttt{zone\_enter}}

Для полигона зоны \(Z\) и траектории \(\tau_j\) используется предикат входа:
\begin{equation}
P_{\text{zone\_enter}} = \mathbb{1}\{p_{t-1}\notin Z \land p_t \in Z\},
\end{equation}
\where
\whereitemlast{\(p_t\)}{реперная точка объекта (центр нижней грани бокса)}

\subsection{Правило \texttt{line\_cross}}

Для ориентированной линии \(L = (a,b)\) вычисляется знак функции ориентации:
\begin{equation}
\sigma_t = \operatorname{sign}\left((b-a)\times(p_t-a)\right).
\end{equation}
Факт пересечения фиксируется при смене знака \(\sigma_{t-1}\sigma_t < 0\) с учетом минимального смещения по времени, чтобы подавить дрожание около линии.

\subsection{Правило \texttt{dwell\_time}}

Для зоны \(Z\) и порога \(\Delta T_{\min}\) правило задается:
\begin{equation}
P_{\text{dwell}} = \mathbb{1}\{p_t \in Z \land (t - t_{\text{enter}}) \ge \Delta T_{\min}\}.
\end{equation}
Событие формируется однократно при первом выполнении условия.

\subsection{Модель \texttt{dedup\_key}}

Для подавления повторной генерации вводится функция
\begin{equation}
\kappa(e) = \texttt{camera\_id} \,\Vert\, \texttt{rule\_id} \,\Vert\, \texttt{track\_id} \,\Vert\, \texttt{bucket}(t),
\end{equation}
\where
\whereitemlast{\(\texttt{bucket}(t)\)}{опциональный временной бакет для правил длительности}

Если ключ \(\kappa(e)\) уже присутствует в скользящем окне дедупликации, событие в таблицу исходящих сообщений не добавляется. Формально:
\begin{equation}
\text{emit}(e) = \mathbb{1}\{\kappa(e)\notin \mathcal{K}_{\text{window}}\}.
\end{equation}
\section{Модель надежной доставки: исходящие сообщения, повторы и очередь недоставленных сообщений}

\subsection{Транзакционная фиксация события}

Пусть транзакция \(T\) вставляет записи в таблицы \texttt{events} и \texttt{event\_outbox}. Гарантия атомарности:
\begin{equation}
\begin{aligned}
\operatorname{commit}(T) \Rightarrow {}&
(\exists e \in \texttt{events})\ \land \\
&(\exists o \in \texttt{event\_outbox}: o.\texttt{event\_id}=e.\texttt{event\_id}).
\end{aligned}
\end{equation}
При откате транзакции (rollback) обе записи отсутствуют \cite{transactional_outbox,richardson_microservices}.

\subsection{Повторные попытки публикации}

Для записи таблицы исходящих сообщений со счетчиком попыток \(n\) время следующей отправки:
\begin{equation}
t_{\text{next}} = t_0 + \min\left(t_{\max}, t_{\text{base}}\cdot 2^n\right).
\end{equation}
Политика ограничивается \(n \le 10\), после чего запись переводится в конечное состояние ошибки.

\subsection{Модель потребителя и очереди недоставленных сообщений}

Интеграционный потребитель очереди HTTP-доставки обрабатывает сообщение и отправляет брокеру подтверждение обработки (ACK) только после успешной доставки HTTP-получателю. При исчерпании лимита попыток сообщение маршрутизируется в очередь недоставленных сообщений \cite{rabbitmq_confirms,rabbitmq_dlq}.

Для прикладной надежности используется приближенная оценка вероятности успешной доставки:
\begin{equation}
P_{\text{succ}} = 1 - \prod_{k=1}^{n}(1-p_k),
\end{equation}
\where
\whereitemlast{\(p_k\)}{вероятность успешной доставки на \(k\)-й попытке}
Рост \(n\) увеличивает \(P_{\text{succ}}\), но повышает задержку и нагрузку на внешний контур.

\section{Оценка вычислительной сложности и ресурсной нагрузки}

\subsection{Сложность стадий конвейера}

Для практического применения важно оценить не только формальную корректность модели, но и ее вычислительную стоимость. Введем обозначения:
\begin{itemize}
    \item \(F_i\) - входной FPS камеры \(i\);
    \item \(K_t\) - число детекций в кадре \(t\);
    \item \(M_t\) - число активных траекторий;
    \item \(R\) - число правил для камеры.
\end{itemize}

Стадия захвата и декодирования имеет стоимость, пропорциональную числу декодируемых кадров:
\begin{equation}
\mathcal{O}_{\text{capture}} \sim \mathcal{O}(F_i).
\end{equation}
При политике замещения последним актуальным кадром часть кадров отбрасывается до инференса, что снижает нагрузку на последующие этапы конвейера.

Стадия детекции имеет наибольшую вычислительную стоимость и определяется архитектурой модели и входным разрешением. В инженерной оценке ее удобно учитывать через измеряемую задержку \(\Delta t_{\text{infer}}\), а не через аналитическую асимптотику.

Для сопровождения ключевым является шаг ассоциации «детекции--треки». При использовании венгерского алгоритма верхняя оценка:
\begin{equation}
\mathcal{O}_{\text{assoc}} \sim \mathcal{O}(\max(K_t, M_t)^3).
\end{equation}
На практике после гейтинга и фильтрации матрица стоимости разрежена, поэтому фактическая стоимость ниже.

Стоимость правилового движка в базовой реализации:
\begin{equation}
\mathcal{O}_{\text{rules}} \sim \mathcal{O}(M_t \cdot R),
\end{equation}
поскольку каждая активная траектория проверяется по набору правил камеры.

\subsection{Задержка сквозного контура}

Суммарная задержка от объекта к событию представляется как
\begin{equation}
\Delta t_{\text{obj}\to\text{event}} =
\Delta t_{\text{capture}} +
\Delta t_{\text{queue}} +
\Delta t_{\text{infer}} +
\Delta t_{\text{track}} +
\Delta t_{\text{rules}} +
\Delta t_{\text{store}}.
\end{equation}
Если очередь между захватом и инференсом не ограничена, \(\Delta t_{\text{queue}}\) может расти неограниченно. В предлагаемой модели размер очереди фиксирован единицей, поэтому:
\begin{equation}
\Delta t_{\text{queue}} \le \frac{1}{F_i},
\end{equation}
что гарантирует отсутствие лавинообразного накопления задержки.

\subsection{Ресурсные ограничения на одном узле}

Односерверная модель требует явного баланса между числом камер и доступными ресурсами:
\begin{itemize}
    \item центральный процессор (CPU): декодирование, сопровождение объектов, правила, сериализация событий;
    \item аппаратный ускоритель: инференс детектора при соответствующем стенде;
    \item оперативная память (RAM): буферы кадров, состояния траекторий, кэш конфигурации;
    \item дисковый ввод-вывод: запись событий, таблицы исходящих сообщений и артефактов.
\end{itemize}

Для камеры \(i\) можно использовать приближенную модель использования памяти:
\begin{equation}
Mem_i \approx Mem_{\text{frame}} + Mem_{\text{tracks}} + Mem_{\text{exec}},
\end{equation}
\where
\whereitem{\(Mem_{\text{frame}}\)}{память, занимаемая кадровым буфером}
\whereitem{\(Mem_{\text{tracks}}\)}{память состояния активных траекторий}
\whereitemlast{\(Mem_{\text{exec}}\)}{накладные расходы модели и библиотек в контуре исполнения}

Эта оценка используется для первичной оценки емкости стенда перед приемочными прогонами.

\section{Псевдокод целевого алгоритма}

\subsection{Алгоритм обработки одной камеры}

Ниже приведен укрупненный алгоритм контура исполнения для камеры \(i\):

\begin{enumerate}
    \item Инициализировать источник потока, модуль сопровождения и локальный кэш правил.
    \item Пока камера активна:
    \begin{enumerate}
        \item Получить очередной кадр из RTSP/MP4.
        \item Нормализовать временную метку и записать кадр в буфер последнего актуального кадра (замещая предыдущий).
        \item Если есть актуальный кадр для инференса, запустить детектор.
        \item Отфильтровать детекции по порогам достоверности и IoU, а также по целевым классам.
        \item Обновить модуль сопровождения и получить список активных траекторий.
        \item Проверить правила \texttt{zone\_enter}, \texttt{line\_cross}, \texttt{dwell\_time}.
        \item Для каждого нового факта сформировать \texttt{dedup\_key}.
        \item Если ключ не встречался в окне дедупликации, в одной транзакции записать:
        \begin{enumerate}
            \item запись в \texttt{events};
            \item запись в \texttt{event\_outbox}.
        \end{enumerate}
    \end{enumerate}
\end{enumerate}

Такое представление соответствует архитектурной декомпозиции из главы 3 и позволяет связать математическую модель с кодовой реализацией модулей.

\subsection{Алгоритм работы публикатора исходящих сообщений}

\begin{enumerate}
    \item Выбрать ограниченную партию записей таблицы исходящих сообщений со статусом \texttt{pending/retry}.
    \item Для каждой записи:
    \begin{enumerate}
        \item Сформировать ключ маршрутизации вида \texttt{event.<type>.v1}.
        \item Опубликовать сообщение в обменник \texttt{events.topic.v1}.
        \item При подтверждении публикации отметить запись как \texttt{published}.
        \item При ошибке увеличить \texttt{retry\_count} и вычислить \texttt{next\_retry\_at} по политике задержки между повторами.
    \end{enumerate}
    \item Повторять цикл с малым интервалом опроса.
\end{enumerate}

\subsection{Алгоритм доставки HTTP-получателю}

\begin{enumerate}
    \item Получить сообщение из очереди \texttt{integration.webhook.v1}.
    \item Выполнить HTTP POST во внешний конечный адрес.
    \item Зафиксировать попытку в \texttt{delivery\_attempts}.
        \item При успехе отправить брокеру подтверждение обработки (ACK).
    \item При ошибке:
    \begin{enumerate}
        \item если лимит попыток не исчерпан, инициировать повтор;
        \item иначе отправить сообщение в очередь недоставленных сообщений и пометить конечную ошибку доставки.
    \end{enumerate}
\end{enumerate}

\section{Параметризация и калибровка модели}

\subsection{Набор настраиваемых параметров}

Для воспроизводимых экспериментов фиксируется конфигурационный вектор:
\begin{equation}
\Theta = \{\tau_{\text{conf}}, \tau_{\text{iou}}, \lambda, T_{\text{lost}}, \Delta T_{\min}, n_{\max}, t_{\text{base}}, t_{\max}\}.
\end{equation}
\where
\whereitem{\(\tau_{\text{conf}}\)}{порог достоверности детекции}
\whereitem{\(\tau_{\text{iou}}\)}{порог IoU для NMS}
\whereitem{\(\lambda\)}{баланс признаков движения и внешнего вида в ассоциации}
\whereitem{\(T_{\text{lost}}\)}{таймаут удержания потерянного трека}
\whereitem{\(\Delta T_{\min}\)}{порог для \texttt{dwell\_time}}
\whereitem{\(n_{\max}\)}{максимальное число повторных попыток}
\whereitem{\(t_{\text{base}}\)}{базовый интервал между повторами}
\whereitemlast{\(t_{\max}\)}{максимальный интервал между повторами}

\subsection{Принцип калибровки}

Калибровка выполняется поэтапно:
\begin{enumerate}
    \item подобрать \(\tau_{\text{conf}}, \tau_{\text{iou}}\) на валидационном подмножестве кадров;
    \item зафиксировать параметры модуля сопровождения (\(\lambda, T_{\text{lost}}\)) по критерию минимизации \(ID\)-switches;
    \item настроить пороги правил так, чтобы снизить ложные события без роста пропусков;
    \item выбрать параметры повторной доставки и задержки между повторами по компромиссу «надежность--задержка».
\end{enumerate}

Важный практический момент: оптимум по одной метрике (например, mAP детектора) не гарантирует оптимум по системной метрике \(\operatorname{p95}_{\text{obj}\to\text{event}}\). Поэтому окончательная калибровка должна проводиться в целевом видеоконтуре, а текущая работа фиксирует модель такой калибровки и границы стендовой проверки.

\subsection{Ограничения валидности модели}

Текущая формализация имеет осознанные ограничения:
\begin{itemize}
    \item нет межкамерной реидентификации;
    \item правиловой движок ориентирован на атомарные события и не покрывает сложные составные сценарии;
    \item оценки производительности чувствительны к характеристикам видеопотока и стенда.
\end{itemize}

Несмотря на это, модель достаточна для целей текущего исследования и для проведения корректной экспериментальной оценки, предусмотренной главой 4.

\section{Метрики качества подсистем и их агрегация}

\subsection{Метрики детекции}

Для детектора в offline-контуре фиксируются:
\begin{itemize}
    \item \(AP_{50}\), \(AP_{75}\);
    \item \(mAP_{50:95}\);
    \item precision/recall по целевым классам.
\end{itemize}

При переносе в online-контур детектор оценивается не только по точности, но и по \(\Delta t_{\text{infer}}\), поскольку именно эта величина напрямую влияет на \(\operatorname{p95}_{\text{obj}\to\text{event}}\).

\subsection{Метрики сопровождения}

Для модуля сопровождения используются метрики:
\begin{itemize}
    \item \(MOTA\) - агрегированная точность сопровождения;
    \item \(MOTP\) - точность локализации в сопровождении;
    \item \(IDF1\) - качество сохранения идентичности;
    \item \(ID\text{-}switches\) - число переключений ID.
\end{itemize}

Для задач охранного мониторинга ключевыми являются \(IDF1\) и \(ID\text{-}switches\), поскольку они определяют корректность событий, зависящих от длительности и направления движения объекта.

\subsection{Метрики событийного уровня}

Событийный модуль оценивается как задача бинарной классификации фактов:
\begin{equation}
\text{precision}_e = \frac{TP_e}{TP_e + FP_e}, \quad
\text{recall}_e = \frac{TP_e}{TP_e + FN_e}.
\end{equation}
Дополнительно фиксируется доля подавленных дубликатов:
\begin{equation}
R_{\text{dedup}} = \frac{N_{\text{suppressed}}}{N_{\text{raw\_triggers}}}.
\end{equation}
Слишком низкое \(R_{\text{dedup}}\) может означать недоиспользование механизма дедупликации, слишком высокое - риск агрессивного подавления валидных событий.

\subsection{Метрики надежности интеграции}

Для контура доставки:
\begin{equation}
R_{\text{publish}} = \frac{N_{\text{published}}}{N_{\text{outbox}}}, \quad
R_{\text{webhook}} = \frac{N_{\text{webhook\_ok}}}{N_{\text{published}}}.
\end{equation}
Также учитываются:
\begin{itemize}
    \item среднее число повторных попыток на событие;
    \item доля сообщений, переведенных в очередь недоставленных сообщений;
    \item время обработки накопленной очереди исходящих сообщений после восстановления брокера.
\end{itemize}

\subsection{Сквозная агрегированная оценка}

Для интегральной оценки качества платформы вводится вектор
\begin{equation}
\mathbf{Q} = \left[
\operatorname{p95}_{\text{obj}\to\text{event}},
IDF1,
\text{precision}_e,
R_{\text{webhook}},
R_{\text{drop}}
\right].
\end{equation}
Решение считается приемлемым, если каждая компонента \(\mathbf{Q}\) удовлетворяет целевому порогу. Такой подход предпочтительнее единой скалярной функции, поскольку позволяет локализовать конкретный деградирующий слой конвейера.

\section{Формализация ограничений и критериев приемки}

\subsection{Ограничения по времени}

Пороговые значения в работе задаются как инженерные критерии стендовой приемки: они не подменяют сертификационные процедуры, но согласуются с логикой измерения производительности интеллектуальной видеоаналитики в системах видеонаблюдения~\cite{iec_62676_6_2026}. Для CPU-профиля:
\begin{equation}
\operatorname{p95}_{\text{obj}\to\text{event}} \le 2.0\ \text{с}.
\end{equation}
Для расширенных многокамерных профилей пороги должны устанавливаться только после отдельной серии измерений на целевом оборудовании. Нарушение выбранных ограничений требует пересмотра параметров \(\Theta\) или масштабирования вычислительных ресурсов.

\subsection{Ограничения по надежности}

Требование надежности формализуется как:
\begin{equation}
N_{\text{lost}} = 0 \quad \text{при кратковременной недоступности RabbitMQ},
\end{equation}
а также
\begin{equation}
\frac{N_{\text{dlq}}}{N_{\text{generated}}} \le \epsilon_{\text{dlq}},
\end{equation}
\where
\whereitemlast{\(\epsilon_{\text{dlq}}\)}{допустимый порог доли сообщений, переведенных в очередь недоставленных сообщений}

\subsection{Ограничения по безопасности}

Для защищенных операций API должна выполняться аксиома разграничения доступа, согласованная с принципами минимально необходимых привилегий, управления доступом и подотчетности действий~\cite{nist_sp800_53r5}:
\begin{equation}
\forall u \in \mathcal{U},\ \forall a \in \mathcal{A}_{\text{protected}}:
\quad permit(u,a) \Leftrightarrow role(u) \in \mathcal{R}(a).
\end{equation}
Нарушение данной аксиомы трактуется как критическая дефектность контура безопасности и блокирует приемку релиза.

\subsection{Критерий архитектурной приемки}

Финальный критерий приемки задается множеством ограничений:
\begin{equation}
\mathcal{C} = \{C_{\text{time}}, C_{\text{reliability}}, C_{\text{security}}, C_{\text{observability}}\}.
\end{equation}
Платформа считается прошедшей архитектурную приемку, если:
\begin{equation}
\bigwedge_{c \in \mathcal{C}} c = 1.
\end{equation}
\where
\whereitem{\(C_{\text{time}}\)}{соблюдены ограничения по задержке}
\whereitem{\(C_{\text{reliability}}\)}{подтверждена корректность таблицы исходящих сообщений, повторов и очереди недоставленных сообщений}
\whereitem{\(C_{\text{security}}\)}{RBAC-ограничения подтверждены тестами}
\whereitemlast{\(C_{\text{observability}}\)}{все обязательные метрики доступны}

Данный критерий используется как формальная основа интерпретации результатов в главе 4.

\section*{Выводы по главе 2}
\addcontentsline{toc}{section}{Выводы по главе 2}

Во второй главе сформированы теоретические предпосылки, модели и методология построения платформы видеоаналитики реального времени. Основные результаты главы могут быть резюмированы тремя положениями.

\textit{Во-первых}, обоснован выбор архитектурного шаблона и технологического стека. Микросервисная архитектура с разделением на контур управления и контур обработки данных согласована с архитектурными факторами, сформулированными в главе 1 (изоляция отказов, локальность изменений, наблюдаемость). Для каждого слоя стека (детектор, модуль сопровождения, веб-фреймворк, СУБД, брокер сообщений, видеозахват, наблюдаемость, оркестрация) выполнено сравнение альтернатив и зафиксирован обоснованный выбор: YOLO~\cite{redmon2016yolo,redmon2018yolov3,ultralytics_docs} и DeepSORT~\cite{wojke2017deepsort,kalman1960,kuhn1955hungarian} для алгоритмического ядра, FastAPI~\cite{fastapi_docs} и SQLAlchemy~\cite{sqlalchemy_docs} для управляющего сервиса, PostgreSQL~\cite{postgresql_docs} и RabbitMQ~\cite{rabbitmq_docs,rabbitmq_confirms,rabbitmq_dlq} для надежной фиксации и асинхронной доставки событий, Prometheus и Grafana~\cite{prometheus_docs,grafana_docs} для наблюдаемости, Docker Compose~\cite{docker_docs} для воспроизводимого стенда. Сводный технологический стек отражен в таблице~\ref{tab:stack_summary}, его архитектурное представление - на рисунках~\ref{fig:ch2_pipeline} и~\ref{fig:ch2_c4_container}.

\textit{Во-вторых}, сформирована математическая модель сквозного конвейера платформы: от захвата видеопотока и детекции до сопровождения, правиловой генерации событий и надежной доставки через транзакционную таблицу исходящих сообщений~\cite{transactional_outbox,richardson_microservices}. Зафиксированы формальные определения для трех базовых правил (\texttt{zone\_enter}, \texttt{line\_cross}, \texttt{dwell\_time}), введена модель дедупликации на основе \texttt{dedup\_key}, описаны политика повторной доставки с экспоненциальной задержкой между попытками и маршрутизация в очередь недоставленных сообщений. Диаграмма последовательности доставки события (рисунок~\ref{fig:ch2_seq_outbox}) связывает этапы транзакционной фиксации и асинхронной публикации с конкретными участниками контура.

\textit{В-третьих}, введены оценки вычислительной сложности этапов конвейера и формализован критерий архитектурной приемки $\bigwedge_{c \in \mathcal{C}} c = 1$, объединяющий ограничения по времени, надежности, безопасности и наблюдаемости. Зафиксированы наборы метрик подсистем (детекции, сопровождения, событийного уровня, надежности интеграции) и сводный вектор интегральной оценки $\mathbf{Q}$. Полученные формализации используются как методическая основа инженерной реализации платформы (глава~\ref{ch:implementation}) и как набор критериев для интерпретации результатов; в главе~\ref{ch:approbation} фактические измерения отделяются от целевых критериев расширенной приемки.
